\documentclass[10pt]{article}
\usepackage[T1]{fontenc}
\usepackage[utf8]{inputenc}
\usepackage[french]{babel}
\usepackage{amsmath,amssymb,mathrsfs,mathpazo}
\usepackage{geometry}
\geometry{a4paper,margin=2cm}
\usepackage{graphicx,array,tabularx,tabularray,colortbl,multirow,multicol,booktabs,enumitem,pifont,ulem,fancybox,diagbox,cancel}
\usepackage{pstricks,pstricks-add,pst-plot,pst-tree,pst-node,pst-text,pst-func,pst-math,pst-bspline,pst-3dplot,pst-eucl,pst-solides3d}
\let\pscalculate\undefined
\usepackage{pst-bezier}
\usepackage{tikz,pgfplots,tkz-tab}
\usetikzlibrary{arrows,arrows.meta,patterns,calc,positioning,trees,decorations.markings,shapes.geometric,shapes.misc,shapes.arrows,intersections,decorations.pathmorphing,decorations.pathreplacing,backgrounds,shadows,babel}
\pgfplotsset{compat=1.18}
\usepgfplotslibrary{fillbetween}
\usepackage[np]{numprint}
\usepackage{listings,makecell,pas-tableur,twoopt}
\let\barre\undefined
\usepackage[pstricks]{tablvar}
\definecolor{darkgreen}{RGB}{0,100,0}
\providecommand{\red}{\color{red}}
\providecommand{\blue}{\color{blue}}
\providecommand{\green}{\color{green}}
\providecommand{\black}{\color{black}}
\providecommand{\white}{\color{white}}
\providecommand{\gray}{\color{gray}}
\providecommand{\cyan}{\color{cyan}}
\providecommand{\magenta}{\color{magenta}}
\providecommand{\yellow}{\color{yellow}}
\providecommand{\violet}{\color{violet}}
\providecommand{\orange}{\color{orange}}
\providecommand{\euro}{\texteuro}
\providecommand{\EUR}{\texteuro}
\providecommand{\ieme}{\textsuperscript{e}}
\providecommand{\degres}{\ensuremath{^\circ}}
\providecommand{\degree}{\ensuremath{^\circ}}
\providecommand{\vv}[1]{\overrightarrow{#1}}
\providecommand{\Coord}[1]{#1}
\providecommand{\decosix}{}
\providecommand{\starredbullet}{\ensuremath{\bullet}}
\pagestyle{empty}
\setlength{\parindent}{0pt}
\usepackage{adjustbox}\DeclareUnicodeCharacter{2236}{:}\DeclareUnicodeCharacter{2190}{\ensuremath{\leftarrow}}
\providecommand{\R}{\mathbb{R}}
\providecommand{\N}{\mathbb{N}}
\providecommand{\Z}{\mathbb{Z}}
\providecommand{\Q}{\mathbb{Q}}
\providecommand{\D}{\mathbb{D}}
\providecommand{\C}{\mathbb{C}}
\providecommand{\e}{{\mathrm{e}}}
\providecommand{\vect}[1]{\overrightarrow{#1}}
\providecommand{\vectt}[1]{\overrightarrow{\mathrm{#1}}}
\providecommand{\vvt}[1]{\overrightarrow{\mathrm{#1}}}
\providecommand{\np}[2][]{{#2~#1}}
\providecommand{\vv}[1]{{\overrightarrow{#1}}}
\renewcommand{\d}{{\mathrm{d}}}
\providecommand{\degres}{{^\circ}}
\providecommand{\degre}{{^\circ}}
\providecommand{\blue}{\color{blue}}
\providecommand{\red}{\color{red}}
\providecommand{\enskip}{\quad}
\providecommand{\llbracket}{\lbrack\!\lbrack}
\providecommand{\rrbracket}{\rbrack\!\rbrack}
\providecommand{\euro}{{\text{€}}}
\providecommand{\no}{{\text{n°}}}
\providecommand{\interval}[2]{{[#1;#2]}}
\providecommand{\intervFO}[2]{{[#1;#2[}}
\providecommand{\dotfill}{\ldots}

\providecommand{\ds}{\displaystyle}
\providecommand{\textding}[1]{#1}
\providecommand{\Oij}{\(\mathrm{O};\vec{\imath},\vec{\jmath}\)}
\providecommand{\Oijk}{\(\mathrm{O};\vec{\imath},\vec{\jmath},\vec{k}\)}
\providecommand{\Ouv}{\(\mathrm{O};\vec{u},\vec{v}\)}
\providecommand{\intervFF}[2]{\ensuremath{[#1;#2]}}
\renewcommand{\euro}{\ifmmode\text{\texteuro}\else\texteuro\fi}
\DeclareUnicodeCharacter{00B0}{\ifmmode^{\circ}\else\textdegree\fi}
\DeclareUnicodeCharacter{20AC}{\ifmmode\text{\texteuro}\else\texteuro\fi}
% Énoncés : Éducation nationale. Transcription et corrigés : APMEP, Denis Vergès et contributeurs. Découpage et classement : CoopMaths, Sébastien Lozano, Éric Elter et contributeurs. Adaptation de la présentation : Maths Explorer.
% CC BY-SA : https://coopmaths.fr/www/about/
% https://coopmaths.fr/alea/static/eam/2025/tex/eam_specifique_2025_06_sujet0v3_2.tex
\begin{document}
\noindent\textit{APMEP / CoopMaths. CC BY-SA. Présentation adaptée par Maths Explorer.}\par\medskip
Victor sort un plat du four. La température du plat est alors égale à 180\ensuremath{^\circ} C. Il place ce plat dans une pièce dont la température est égale à 25\ensuremath{^\circ} C. Le plat refroidit. \\
Le plat ne pourra être servi que lorsque sa température sera devenue inférieure ou égale à 40\ensuremath{^\circ} C.

On étudie le refroidissement du plat selon deux modèles mathématiques.

\subsubsection*{Partie A : Premier modèle.}
On suppose que la baisse de la température du plat est \textit{proportionnelle} à la durée du refroidissement, c'est-à-dire au nombre de minutes écoulées depuis la sortie du four.

On constate que 3 minutes après la sortie du four, la température du plat est égale à 105\ensuremath{^\circ} C.

\begin{enumerate}
\item De combien de degrés le plat a-t-il baissé en 3 minutes ? En 1 minute ?

\item Vérifier que la température du plat, 5 minutes après la sortie du four, est égale à 55\ensuremath{^\circ} C.

\item Selon ce modèle, quelle serait la température du plat, 8 minutes après la sortie du four ? Ce premier modèle semble-t-il pertinent ?
\end{enumerate}

\subsubsection*{Partie B : Second modèle.}
On dispose toujours des données suivantes :
\begin{itemize}
\item la température de la pièce est égale à 25\ensuremath{^\circ} C
\item la température du plat à la sortie du four est égale à 180\ensuremath{^\circ} C
\item la température du plat, 3 minutes après la sortie du four, est égale à 105\ensuremath{^\circ} C
\end{itemize}

Pour tout entier naturel $n$ on note $U_n$, la différence entre la température du plat et la température de la pièce, $n$ minutes après la sortie du four.

\underline{Exemple :} 3 minutes après la sortie du four, l'écart avec la température de la pièce est égal à $105 - 25 = 80$. On a donc $U_3 = 80$.

\begin{enumerate}
\item Justifier que $U_0 = 155$.

\item On suppose que chaque minute la différence $U_n$ diminue de 20\%.
\begin{enumerate}
\item Justifier que, pour tout entier naturel $n$, on a $U_{n+1} = 0,8U_n$.
\item En déduire la nature de la suite $(U_n)$ et donner sa raison.
\item Exprimer $U_n$ en fonction de $n$, pour tout entier naturel $n$.
\item On dispose des données suivantes :
\end{enumerate}
\end{enumerate}

\vspace{-0.9\baselineskip}
\begin{center}
\begin{tabular}{|c|c|c|c|c|c|c|c|c|c|c|c|c|c|}
\hline
$n$ & 3 & 4 & 5 & 6 & 7 & 8 & 9 & 10 & 11 & 12 & 13 & 14 & 15 \\
\hline
$U_n$ & 80 & 64 & 51,2 & 41 & 32,8 & 26,2 & 21 & 16,8 & 13,4 & 10,7 & 8,6 & 6,9 & 5,5 \\
\hline
\end{tabular}
\end{center}

\vspace{-0.9\baselineskip}
\[U_n \text{ arrondi à } 10^{-1}\]

\vspace{-0.5\baselineskip}
\begin{itemize}
    \item[]
    \begin{itemize}
    		\item[] Au bout de combien de minutes, Victor pourra-t-il servir le plat ?
	\end{itemize}
\end{itemize}

\end{document}
